\textbf{Identification almost surely.} Fix \(a\in\{0,1\}\). By \(\mathrm{ass:conditional-exchangeability}\),
+\[
+\mathbb E_P[Y^a\mid X,A=a]=\mathbb E_P[Y^a\mid X]
+\]
+at every covariate value where the arm-specific conditional expectation is defined. By \(\mathrm{ass:strict-overlap}\), \(P(A=a\mid X)\ge c_\pi>0\) almost surely, so this identity is available for \(P_X\)-almost every covariate value. On \(\{A=a\}\), \(\mathrm{ass:consistency}\) gives \(Y=Y^a\). Consequently,
+\[
+\mu_a(X)=\mathbb E_P[Y\mid A=a,X]=\mathbb E_P[Y^a\mid A=a,X]=\mathbb E_P[Y^a\mid X]
+\]
+almost surely. Subtracting the identities for \(a=0\) and \(a=1\), using linearity of conditional expectation and \(\mathrm{def:cate}\), yields
+\[
+\tau_P(X)=\mu_1(X)-\mu_0(X)=\mathbb E_P[Y^1-Y^0\mid X]
+\]
+almost surely.
+
+\textbf{Uniqueness of the continuous representative.} Let \(g_1,g_2\) be continuous functions on \(\mathcal X\) representing the same \(P_X\)-almost-sure equivalence class. Suppose that \(g_1(x_0)\ne g_2(x_0)\) for some \(x_0\in\operatorname{supp}(P_X)\). Continuity gives \(\varepsilon>0\) and a relatively open neighborhood \(O\subseteq\mathcal X\) of \(x_0\) such that \(|g_1(x)-g_2(x)|>\varepsilon\) for every \(x\in O\). Since \(x_0\in\operatorname{supp}(P_X)\), one has \(P_X(O)>0\), contradicting equality of \(g_1\) and \(g_2\) almost surely. Thus \(g_1=g_2\) on \(\operatorname{supp}(P_X)\). Applying this argument under \(\mathrm{ass:continuous-cate-version}\) proves uniqueness of the designated continuous CATE representative there.
+
+The arm regressions \(\mu_0\) and \(\mu_1\) are observed-law conditional-mean functionals, and strict overlap determines both of their \(P_X\)-almost-sure classes. Hence \(\mathrm{def:cate}\) and \(\mathrm{def:upper-set}\) make \(G_P\) an observed-law functional modulo \(P_X\)-null sets. Finally, \(\mathrm{ass:bounded-density}\) and \(\mathrm{def:weighted-loss}\) give
+\[
+d_H(H,G_P;P)=\int_{H\triangle G_P}|\tau_P(x)-\theta|f_P(x)\,dx
+=\int_{\mathcal X}\mathbf 1_{H\triangle G_P}(x)|\tau_P(x)-\theta|\,dP_X(x).
+\]
+Changing either \(H\) or \(G_P\) on a \(P_X\)-null set changes the last integrand only on a \(P_X\)-null set and therefore leaves the integral unchanged. This proves the asserted invariance.